\section{The Esoteric Origin of the Species}

\begin{quotex}
This essay originally appeared in the \emph{Introduction to Magic}, volume 3. It was published under the name ``Avro'', which I believe to be \textbf{Julius Evola}.
\end{quotex}

\paragraph{Part I}
The reader will not fail to notice that in one particular point among others, the esoteric teachings have a distinctly counter-current character: this is in regards to ``evolution''.

Evolution is a type of obsession of the modern mind. It is a true psychological complex that, in the middle of the ``logic of the subsoil'' about which Iagla Evola wrote (Vol 2, ch 2), directs the minds of many of those who presume that they were using the ``scientific method'' and objective research. Here it would be necessary to be convinced of what is also of value for many other things: or else that certain possibilities of understanding, seeing, and controlling are the effect of a certain change of attitude: and not vice versa, as rationalism would like.

Regarding evolution, it could, for example, have caused surprise about what we said about the Hyperborean tradition. In effect, the idea that previously in the interglacial and paleolithic period, a great unitary civilization existed, to which the fundamental symbols, the roots of language, and the graphism\footnote{\url{https://en.wikipedia.org/wiki/Graphism}} are related to the oldest cultures --- a similar idea must arise as revolutionary in respect to modern opinions, which are believed to have been proved once and for all on a positive basis. And it is not only a question of simple evolutionism in the history of civilization: even parts of science begin to be touched where, in one form or another, the Darwinian hypothesis on the origins of the species and on the descent of man from animals is current even now.

So the problem must be confronted as a whole. We will therefore mention what the esoteric teaching say about it, without going into details that could take me rather far: beyond the extent it can be usefully discussed in this journal.

Although not stopping myself here, I have to emphasize first of all that in the same field of profane science today it is no longer about the evolutionism that Darwin announced in his time. The Darwinian hypotheses, ever since, has undergone numerous criticisms, and real difficulties, previously not seen, have compelled it to be modified.

Where it is especially shown to be weak, is in its attempt to deduce the variety of the species from an almost automatic play of material conditions of the environment, of natural selection, and of inherited transmission of acquired characteristics. The vitalistic point of view, which asserts precedence and an excess of vital energy over all the conditions, must instead acquire a greater significance.

\textbf{Bergson} is among those who have come into play in strictly scientific terrains against Darwinian evolutionism, showing through its insufficiencies that have left the field free for the hypothesis of an evolutionism of a no longer bio-materialistic, but of a creative, character.

The same biologists after all recognized how much the variety of species resists the attempt of that simple and linear deduction, the one from the other, that Darwin conjectured; they were led to admit ``quasi-jumps'' from some species to others, so that the hypothesis most in vogue today is that of De Vries, who appeals to unpredictable and essential internal mutations precisely to explain those jumps within the scheme of evolutionism.

This, nevertheless, is one hypothesis like any other, which is of interest only for the difficulty that it calls attention to, and that in a great part exists even after the introduction of enigmatic concept of mutation. I note in passing that the difficulty finds an exact comparison in what physics has recently chanced upon, when, with the theory of quanta (Planck) and the principle of indeterminism (Heisenberg), it has had to stop in the face of finite quantities of ``action'', without being able to explain, and without being able to point out a continuous process that leads from one to the other.

Moreover, since the evolutionistic hypothesis continues to be in force in biology as far as it has been modified and revised, it is necessary to examine its general foundations. We know the impressive bounty of facts Darwin and his school have gathered in the field of morphology, embryology, paleontology, and even geology. No one would consider denying these facts. What there is to discuss, and to indicate as arbitrary or, at least, as one-sided, is in the interpretation, through which in Darwinism these facts become proofs to support the materialist evolutionistic concept. Without postponing any longer, I will go directly to the fundamental argument.

\emph{When they had just succeeded in establishing a continuity of forms and rings, permitting the passage from one species to another, up to man, with that, one line remains simply established, and no one tells us in what direction it has traversed.} Hence, a priori, every fact adopted to uphold evolutionism could be simultaneously adopted to sustain the opposite thesis: the \emph{involutionistic} thesis: neither more nor less. As true that the lower species are the levels preceding the higher---it is as true that they are instead \emph{degenerative involutions} of the latter. \emph{The presence of intermediate stages} (when they are also transitional stages)---and not from \emph{crossbreeding} or from \emph{selection}: other possibilities that the evolutionists don't consider) \emph{cannot tell me, by themselves, in what direction the march happened}.

That is the fundamental point. Let's now see what can be added to it.

Let us begin with the peoples called backward from the point of view of their mentality and their civilization. Who tells us that they represent the primitive states of present humanity rather than involutive forms and residues of an even older humanity? The fact that the backward peoples are more likely to disappear than to ``evolve'', should make us reflect. Besides, one must consider that an even more ancient humanity could have been \emph{different}, as much from not leaving traces where forms of civilization closer to us were successful and superimposed: from leaving instead traces of them precisely in their degenerating origins, but only of the same trunk. Modern ethnological studies on presumed ``primitives'' have ascertained in them not a lower level of the same mentality, but rather of another mentality, another civilization. From them, through integration, one can reach back to an even older humanity. Modern explorers of prehistory, like Frobenius and Wirth, for example, have followed this method.

\paragraph{Part II}

Moving from primitive man to the anthropoid and the ape, and presuming the jumps are surmountable, which therefore must be made in order to reach the other animal forms according to transformist models, we can say the same thing: as far as being able to consider many animal species as degenerations or degradations of still older non-animal forms. Our point of view is exactly this: \emph{Man does not derive from the animal but, if ever, it is the various animal species that derive from man, in a sense that we now try to explain.}

The principle difficulty that this point of view encounters lies in the fact that the traces of man end at a given geological period: while the traces of prehistoric animals continue up to a much older period. But this same fact is susceptible to different interpretations, for those who can consider the idea of transformations with sufficient breadth: \emph{that the mineral traces of man are more recent, could mean only that man was the last to come into that process, under a certain involutional aspect, through which it is precisely possible that they exist as fossils and therefore those traces are retrievable.}

The misunderstanding about ``cavemen'' comes from not thinking that it is natural that certain very old traces are still found in the caves, because, for multiple factors, they have not been preserved elsewhere. The idea of the recent appearance of man on the earth is based on a lack of vision of the same type; I certainly am not asking anyone to admit the fall of man from the sky: it is enough to surpass the \emph{mineral}, but not the material, concept of corporality; it suffices to think of the possibility of a body whose most physical element (which today is the skeletal system) was composed by a substance unsusceptible to self-conservation through the process of fossilization---to remain indifferent in the face of the fact of the lack of traces in the most remote geologic periods, and to be able to admit the existence of primordial human stocks (of which the anthropoids would be the first degenerating materializations) coexisting among more driven forms of the involutive process, which would be represented by the oldest animals of prehistory. This concept has nothing absurd in itself. Analogically, every manifestation has of necessity an inversive character: whatever stands more to the origin, more to the internal, more to the center, can only be the last to appear in the movement toward the external. And at the center and at the origin, according to esoteric teaching, would exactly be Man.

NOTE: What happens to the rest in every finalistic process: the goal, the end precedes as idea all those conditions that are necessary to actualize it, and as reality it appears as the last, after them.

This Man naturally does not correspond to the man of today: but he corresponds to them in the sense that man of today can be considered as the most approximate manifestation and the most direct lineage of primordial Man. As such, he represents the origin, the axis, while the other species represent lateral or divergent directions, if not to say, byproducts.

With an image previously used by ``Ea'', I can perhaps indicate better what otherwise would require a long detour of concepts in order to be understood.

Let us imagine something like an assault, like an undertaking of conquest. A group of tightly unified forces confront the danger, they approach the goal. A battle is undertaken. On the defenseless front some begin to fall back, others advance. The resistance is met, the clash begins. The empire bends. Few succeed in maintaining the original direction -- like the wake of a boat, they leave behind them, dispersed, stopped, sacrificed or beaten, the greater part of those, with whom they were together at the beginning of the undertaking. The group of survivors hold firm, advance again: always fighting, they succeeds in opening a way, they finally reach the position that it was the goal of the undertaking, they conquers it, they maintain it, they planted their own banner in it. In all the other directions, the aborted, frustrated, arrested attempts of the same will lag behind on the way without exit.

We mean now by ``conquest'', the achievement of the physical state of existence in the conditions well-known to us today; we mean that those who have reached the end are men and others who were defeated or have deviated, are the various animal species. Man, as known today, expresses the form in which the stock of a primordial humanity is conserved and still standing in the conditions of materiality, in which various animal species were included in the beginning, certainly not in the forms that appear to us today, but rather in the principles of them, which have taken their origin biologically from a degenerative specialization, from a ramification of the original direction into divergent directions, each one expressing the exhaustion of an attempt, the arrest of a surge of assault from which those who have ``broken down'' are separated and that they left behind themselves.

It is interesting to notice that ideas of this type were also put forward at the fringes of modern culture, without a connection with traditional teaching. In the works of \textbf{Edgar Dacqué} Leben als Symbol, Urwelt, Sage und Menschheit I find, for example, a rather similar conception, and an illustration that makes very clear the concept discussed above.

\paragraph{Conclusion}
\begin{quotex}
The Diagram discussed is by \textbf{Edgar Dacqué}. He was a German paleontologist who was a devout Lutheran, a Theosophist, and was involved with the OTO.

In this conclusion, Avro explains how man arises from the collective. He makes use of the Aristotelean-Thomist teaching in an interesting way. The weakness of the A-T system is that it does not take the individual into account. Here Avro expands on that teaching, with the person as the form or idea and the genotype as the material or matter.
\end{quotex}


In U, where it runs in the central direction, there is man as we know him. A, B, etc., can be considered as the nearest attempts, e.g., the diluvial man, the anthropoid, the ape-man. Further down and backward in time, there are the other, more elementary and divergent, attempts, the first to appear in the completely ``densified'' form of existence that we know, accessible to paleontology and whose region extends to the outside of the circular line marked in the figure. The individual points of exit on this line from which then a group of vectors and secondary deviations starts in turn, represent the ``types'' of each species; and these secondary deviations (which are, in the figure, the vectors U'U'', A'A'', B'B'', F'F'', etc., in respect to the directions U, A, B, F, which also continue to the outside the vector that comes from the axis) are the transformation that each species has undergone in a partial battle, in the middle of adaptation, selection, etc.: that is to say, in the middle of the factors to which Darwin wanted to reduce the whole. Instead, the passage from one species to another does not happen at the periphery: at the periphery only hybrids (as in Z and Y) can be verified, erroneously interpreted by evolutionists as ``transitional forms''. The ``passage'' is instead determined by the issuing and from the emerging of a new branch that departs from the central direction (which goes toward man) to the head of a new impulse, at the failure of what preceded.

Here we can also take into account the true meaning of another of the facts that would seem to prove evolutionism, summarized in the idea that ontogeny repeats phylogeny. That man, in his development from the embryo, through a series of phases which have a certain resemblance with the animal forms of life, from our point of view, means only this: that each human realization includes in brief and retraces the attempts, whose possibility was included in the original stock: \emph{but to retrace them exactly on the base of the original impulse, that goes beyond them all}; something evident in itself, insofar as an embryo stopped in any phase you want, always remains just a human embryo, not that of a fish or another animal species.

It is in a very special sense, not material and not valid for biological consideration, that one can say that to the extent the central line UU passes through the origin of the various types of animality, man, before appearing as such in this world, lived in the hierarchy of animal forms. It is not a question, moreover, of any of the animal forms that are terrestrial and historical manifestations anterior to man: instead it is about that which corresponded to the ``sacred animals'' in the ancient mystery traditions. In the cult of such ``animals'', that in an involutive and restricted form often are found again in certain primitive peoples, that conceal the record of this knowledge, of this co-essentiality, being related to other planes or states of existence. The sacred animals are powers of life that made up part of primordial man, and from which he separated to go beyond; they are not individual animals, but the ``group soul'' \emph{anima di gruppo} , daemons, whose body is the total life of a given animal species---and, each one, a transcendent experience, used up and abandoned in the wake of the past through a type of catharsis or purification. But, as I said, they were in origin parts of universal man, and terrestrial man was the most direct expression of the origin the relationships, the hidden correspondence, will remain: those that esoterism considers as marking the limbs, the functions, and the energies of the human body with the zodiac symbols or other equivalents.

Our figure corresponds, more or less, to the model of a branched tree which should be conceived in its external part as a continuous movement of expansion and return, similar to that produced by a heart: that, in order to indicate the appearance of the individuals of each species and their reabsorption into the original stocks to which they belong, and in which the dark will is conserved that continues to be blindly asserted in the dead-end efforts that constitutes it.

And as each one of these efforts constitutes, less a factor of fall and deviation, an approximation of the central will that reaches its realization only in man; equally humanity itself is an endless reiteration of attempts and approximations that achieves a prefect realization of the type only in very rare men. In a certain way, the battle expressed in the preceding image is still open: if man represents, at least on earth, the last conquered position, as the most advanced position, is the most difficult to maintain. The drifters, the fallen, the deserters, those who are overwhelmed by animal passions are again aspirated and redissolved in forces that bring them back, through contacts with the hidden forces of animality, are without number.

The doctrine of metempsychosis comes into play here in its correct light. It is true in the sense that the possibility of an involutive process is real, for which there is not---as in the vulgar interpretation---the passing of the soul of a man into the body of an \emph{individual} animal: but there is instead its reabsorption into the being or archetype that constitutes the form of the animals born of a given species, whose hidden direction is similar to what informed the whole life of such a man.

Recalling what was said many times in these pages, we can moreover think that, also apart from such ``sacrifices'' of men to ``sacred animals'', other possibilities exist and places of reabsorption. In general, the measure in which the human life receives collective influences, likewise indicates that of the reabsorption and the refoundation in the matrix of humanity, the individual serving as the material for other attempts, other blows, that more or less can happen at the center of the target.

In this regard, many things could be said to clarify the esoteric view about conditioned immortality and reincarnation, if only to interpret phenomena correctly, like those of certain ``inheritances'': which---I give it a quick look---do not have for their basis the transmission from individual to individual, but instead a type of ``habit'' picked up from the ``\emph{genius}'' or ``\emph{manes}'' of a given stock: from where every individual who is enucleated in it brings, at the most, a determined characteristic attribute that refers back to a collective influence.

Perhaps I will give myself the opportunity to add something about this. Here, I will conclude by mentioning that the Aristotelian-scholastic terminology, through which the ``genus'' is the \emph{matter} and the individual is the \emph{form}---the more there is individuation, the more there is form and perfection---it is that which best reflects our ideas.

Everything in life that is still collective, shows how much there is in man that is still unaccomplished. Separating himself from all that is not himself, individuating himself absolutely from being only himself, and only himself, man goes beyond the destiny of rebirth, because he accomplishes the goal of that very rebirth. Without further variations and deviations, he then incarnates the pure central vector by ``not having more daemons'' and by doing a single thing on earth with his ``idea'' and his ``Name''.
\flrightit{Posted on 2013-04-14 by Aeneas}

\begin{center}* * *\end{center}

\begin{footnotesize}
\begin{sffamily}
\texttt{Mercurius on 2013-04-16 at 08:14 said:}

Recently, Evola's review of Berdyaev's ``Meaning of History'' was published on Gornahoor; there are likewise some similarities between Evola and Berdyaev's ideas regarding evolution. Bedyaev neither dismisses nor argues against biological data or ``fact''--for him as well, all of the biological and chronological data could very well be accurate, it is the interpretation that differs. Unlike Evola, Berdyaev does not however discuss a ``devolution'', or difference in ``direction''; rather, he stresses that despite bio/morphological changes, teleology is the best argument contra evolution, ``This does not imply the acceptance of man's destiny as the theme of the metaphysics of history. For to admit the manifestation of man's destiny in world history, it is also necessary to admit his preworld existence;that his destiny originated and was determined prior to the establishment of that world of reality where occur all those processes of evolution and development by which evolutionary theory tries to explain both man's origins and his further development''. Somewhere on Gornahoor seemed to be a very similar excerpt from Solovyov, although it is being allusive right now.

\hfill

\texttt{Ammonius420 on 2013-04-17 at 02:04 said:}

How do the Yugas described by Guénon and Hindu tradition relate to geological time? I recall that Fulcanelli wrote a book that in part touched on fossilized lifeforms and their relation to cosmic cycles, has this been found?

\hfill

\texttt{Boreas on 2014-11-19 at 00:47 said:}

Ammonius, regarding the geological record and the Yuga cycles, there is an excellent article at Cakravartin about the more recent past; the second part of the article especially deals with the geological facts in relation to the cycles. Here's the link:

\url{http://www.cakravartin.com/archives/the-end-of-the-kali-yuga-in-2025-unraveling-the-mysteries-of-the-yuga-cycle-by-bibhu-dev-misra}


\hfill

\texttt{taivasvaeltaja on 2023-09-23 at 11:08 said:}

\url{https://taivaanmalja.com/2023/09/23/kaannos-esoteerinen-lajien-alkupera/}

\hfill

\texttt{Mihai on 2013-04-18 at 04:58 said:}

Evola introduces here some complications which are certainly not sustainable by traditional doctrine. From a Christian point of view, at least, Each species manifest a given aspect of the Logos, being, thus, fixed, each having its own reason for existence, separated to the others. Man is the only one who manifests the Image in its entirety.

Transformism is as much a problem when conceived both evolutionary and involutionary, since it introduces a confusion between the various `logoi', making only a difference of degree between one species and another. Nature, however, is a hierarchical order, with essential differences between the four levels of mineral, vegetal, animal and human.

Now if the Image of God in man, no matter the degradation, can never be annulled, and if animals derive from man in an involutionary process, than we are either forced to admit that the animals also have the Image of God within them, or that, somehow, that Image has been along the way annulled.

So this is going to far. It is enough to say that man can degenerate to such a degree that he comes to resemble the animal, but it is only caused by the fact that he reaches the outer limit of the human domain and the features of the inferior one (that of the animals) begin to be reflected on him, especially since he has strayed so far away from the primordial origin. However, between the two domains there remains an impassable limit.

As to the various geological records:

1. They are extremely controversial

2. They are actually quite smaller than the mainstream authorities try to make us think

3. Their `history' is built through a very procroustean method

4. They, themselves, can account for multiple theories.

\hfill

\texttt{Cologero on 2013-04-18 at 12:48 said:}

You bring up interesting points, Mihai, but as the text says, Evola perhaps means by ``derive'' something different from what you express. In the process of involution, the idea of Man comes first. So these ``derivations'' are actually successive attempts to lead up to the appearance of man in his corporeality. Several of the attempts branched off into different directions, or went extinct. One path led to man. Thus, they are not physical human beings who were degraded. There are two questions, then.

(1) Does this theory accord with the facts as we know them?

(2) Does this theory explain how it leads to Man?

Darwinism answers (1) but not (2), since there is no necessary reason for Man to arise within that theory.

The essential differences between the four kingdoms of mineral, vegetable, animal and human are not absolute. The higher stage is related to the lower stage as genus and differentia. So, while man is an ``animal with reason'', it would not be correct to understand an animal as a ``man without reason''; I think that is the point you are getting at? But I don't see Evola denying that.

You may want to go back and read the essay on Plotinus from the UR group that was published on Gornnahoor a few weeks ago, since a similar them arises. Plotinus claimed that evil was ``derived'' through ``degradation'' and ``dissolution''. So there is a convergence of ideas, as one would expect.

\hfill

\texttt{Pickman on 2013-04-20 at 21:55 said:}

So what came first, the chicken or the egg? Certainly in respect to Darwinism ``there is no necessary reason for Man to arise within that theory'' as the lower does not give birth to the higher and we are left to ponder the ``why'' behind the theory -- since basic bio-determinism, within the limits of a materialist theory or survival of the fittest is ill-equipped to answer what is essentially spiritual.

The goal of the universe was to lead the material to it's superior form of spiritual transgression. Yet we are left with questions as to how the original manifestation began (a fall?). This I find the most difficult part to grasp of the involution theory.

So how did man first appear?

I will leave that question open to others before I venture to offer my own theories on this subject.

\hfill

\texttt{Cologero on 2013-04-21 at 11:18 said:}

Yes, Pickman, that would be a good discussion. Perhaps it should wait until the rest of the essay is made available. Perhaps this idea can be pondered in the meantime. When you consider any of your own projects, the idea of it is formulated and a plan adopted. Then there is the effort to materialize it. Here is where the difficulties arise. Certain aspects of the plan may hit roadblocks or veer into dead ends. A ``feedback loop'' is necessary to continuously alter the plan to take extenuating circumstances into account. The ``material'' world, dominated as it is by the principle of Destiny, or determinism, is resistant to anything. new. It is a measure of one's will to power to be able to overcome such obstacles in order to realize one's ideas in the material world; in this case, essence and existence become a unity. Guenon refers to that as the self-manifestation of one's possibilities. So the ``goal'' of the universe is not really so much to lead the material (which really does not exist at all independently) to the spiritual, but rather to materialize all the possibilities of manifestation.

\hfill

\texttt{Sigurd on 2017-07-07 at 11:32 said:}

``essence and existence become a unity. Guenon refers to that as the self-manifestation of one's possibilities. So the ``goal'' of the universe is not really so much to lead the material (which really does not exist at all independently) to the spiritual, but rather to materialize all the possibilities of manifestation.''

Cologero, could you elaborate on this? It appears very nebulous to me. I presume you aren't referring to some ``horizontal'' application of possibilities; lest you fall into some justification to ``realize'' all sorts of abominations, their very possibility demanding it. You're aware of this, I'm sure: can you edify me on what it means to ``realize'' the spiritual? Is this related to the so-called ``resurrection resolution of the flesh''?

If I'm on the right track, I can employ certain terms to ask for clarification as well, at the risk of misusing them due to my nascent understanding of tradition and philosophy: the essence (eidos), that is virtual, becomes actual or ``substantial'' (ousia) through participation in being (on), its ``material'' side (materia secunda) then giving rise to specific accidents in time.

\hfill

\texttt{Cologero on 2017-07-07 at 22:22 said:}

It's interesting to me, Sigurd, that you even noticed that. It is hard to accept, although I'm not offering it as a \emph{justification} for abominations, but simply as an observable fact. I can outline a little syllogism, and then you can decide if it makes any sense, or might require some refinements (if not outright rejection).

1) We can stipulate that in God, essence and existence are identical.

2) We can stipulate that, although man is both a spiritual and material being, his essence and existence are not identical.

3) Rather, man is a privation, i.e., he has possibilities of being that are not actualized.

4) That is clearly, if trivially, true in the growth and development from infancy to adulthood.

5) I tried to get at this in the \textit{Meditation on the washing of the feet}\footnote{\url{http://www.meditationsonthetarot.com/death-and-resurrection}}. This should be more than a philosophical exercise, but needs to be contemplated. (I.e., raised to the level of ``intellectus'' from the ``ratio''.)

6) I think this is in accord with your third paragraph, isn't it? This is unlike some eastern or gnostic systems which attempt to flee the world, rather than to tend and transform it.

7) As for the ``horizontal'' possibilities that may lead to abominations: are they truly my possibilities? They may be infernal possibilities, i.e., ideas that don't arise from my own Being, but rather are imposed from the outside. As such their ``actualization'' represents a privation of my being, not its actualization. This may be true of seemingly ``good'' accidents, such as one's family, culture, race, etc., insofar as they represent a blind collectivity.

On the cosmic scale, it does indeed seem that all possibilities are being manifested, either in a positive sense or most often as a privation. If Guenon is correct that no possibilities, like snowflakes, are ever exactly the same, then things repeat over and over, just with very subtle variations. For example, how many ways can people die or be tortured? How many ways can humans engage in sex? How many variations in food, clothing, etc. are there? Everything just nauseatingly repeats itself with small changes. And people go wild for it, believing in the novelty of their little variation. The world will end when there has been enough.

So, yes, there is participation in being when the essence becomes actual, but privation is non-being, merely appearance. But the appearances are horrifying.

\hfill

\texttt{Caveman on 2017-12-07 at 19:22 said:}

How should the Neanderthal be understood in this anthrogonic framework? This question is perhaps all the more relevant as it is scientifically alleged today that there occurred some interbreeding between Neanderthals and Cro Magnon man, leaving DNA traces in the modern genome which peaks in Caucasians (and is absent in sub-Saharan Africans). Would they be considered the fallen and deteriorated remnants of some pre-Adamic man? Their apparent practice of cultic rites indicates that some intuition of spiritual reality survived in them, which means that they must probably be considered human, albeit a less central manifestation thereof than their Cro Magnon contemporaries.

\hfill

\texttt{J. R. R. Tolkien on 2020-09-17 at 08:23 said:}

As much as I enjoy Evola's writings, I think he is wrong here.

In the first place, he didn't know that the Neanderthal man was not an ``anthropoid ape like figure'', but was fully human, there were many frauds that were only uncovered recently.

Also the many huge skulls and bones that were found, making the ``evolutionist'' hypothesis even more incredible, bones of giants for example.

To finish it off, Michael Behe proved in solid biological grounds, that mutations cannot construct irreducibly complex forms, because they need all parts at once to work. Not only that, he proved that species tend to devolve, that mutations are actually changes in existing genes, not the creation of new ones.

Lucy, the astrolapithecus, and many other fossils, have been shown to be fully apes. There were many other frauds, that are inummerable. Transitory forms have actually been shown to be very problematic, species appear in the geological record suddenly, as he himself noticed.

Now, to say that there has been a material that made the substance of humans and it doesn't fossilize or show in the geological record is to invent something based on imagination, because there is no proof of that, even though we can clearly see in many ways how the modern life is a degeneration in many ways from the profound ancient wisdom. But there have been improvements too, as in medicine, so there is no progress, nor regression depending on the point of view.

Obviously, spirituality is that which is the highest, so in general one could say that there has been a regression.

But that there were simpler forms to appear I don't think can be denied, again, this only shows that the things more steeped in materiality appeared first, while man appeared last being crowned, there are other elements that will never be known to us regarding the giants, the gigantic skulls, the humans that appear to have degenerated from the primordial stock early on, all we can do is to make up hypothesis that will never lead really anywhere, since the events cannot be brought back, maybe there was an age of gods and men, but how will we know? What about the huge animals that appear in the fossil record, why are there trees that live over thousands of years? Humans have no natural defence against the natural world, no instincts, could it be that humans struggled like beasts in order to ``conquer'' mother Earth? The order of things indicates that humans were given from the beginning the means to achieve his final destination, which is a high spirituality that leads to ecstasy and self-mastery. As soon as they appeared, they were given all of these gifts, but they already existed ideally celestially, in the eternal idea, according to the myths of all traditional peoples they were immortals from the beginning.

\hfill

\texttt{J. R. R. Tolkien on 2020-09-17 at 08:36 said:}

Buried Alive: The Startling Untold Story About Neanderthal Man Jack Cuozzo

The author is an orthodontist who measured the dental arch of the Neanderthals in a Paris museum and found that they had all been assembled wrong, so they looked chinless, he was persecuted fiercely after showing the fraud, this myth of the anthropoid neanderthal is still being told this very day.

\hfill

\texttt{Ciência Aristocrata on 2020-09-18 at 06:17 said:}

The great myths recall divine origins, not animalistic origins. But the Bible speaks of lower races that mingled with the human race, which is curious, the Bible mentions a race of giants, and today we find the fossils of these giants. It is not possible to understand a perfect history of the past with horizontal science, because the history of origins is a metahistory, only a vertical science could say anything about this type of history.

We find human fossils huge and inexplicable skulls, of which ufologists say they are aliens.

But this idea of highly technological aliens is part of the modern myth of progress, which has already been demonstrated by Guénon to be fallacious, the modern myth is ingrained in the modern mind.

If there is purpose in the causes, first humans were given the means to reach their purpose, the purpose of humans is to reach the contemplation of truth.

Materialists who deny the finality of the causes claim that they were like beasts that suffered by dying and attempting to dominate mother nature, in a blind, brute way, this is because the highest thing for materialists is the dominion of matter, they don't see that this actually doesn't satisfy the human will.

Indeed, human beings have no natural defenses, unlike birds that when released into nature know how to distinguish good fruits from poisonous fruits, there is nothing in humans that allow them to distinguish the bad from the good things for their survival, they have no protections against the various dangers of nature, no fur, no shells, nothing, if no science was given to them from the start, not only would they not be able to reach their finality, nor would they be able to survive, it is absolutely necessary that the first humans were aware of the means of protection and to reach the end that is proper only to man, which is contemplation.

\hfill

\texttt{Cologero on 2013-04-22 at 08:44 said:}

On page 180 of the American edition of Revolt, the reference to Dacque was dropped and only the reference to The Transformist Illusion by Douglas Dewar was retained. That whole paragraph has some errors, which are even misleading. E.g., the original text does not contain the phrase ``man is not alone''. One would wonder then who else is far from being the product of evolution. I will be providing a revised translation in this space soon.

\hfill

\texttt{Max on 2015-07-24 at 06:22 said:}

Evola says this: ``... humanity itself is an endless reiteration of attempts and approximations that achieves a perfect realization of the type only in very rare men.''

I interpret it as that the race is not necessarily what determines a person, but on the contrary a person possibly determines the race. The other individuals who belong to the same race are ``approximations'' to its fullest expression. From this it follows that there are guiding ideas of races. He writes that races and species are different ``directions''. I think we could characterize the horizontal direction in a race, species or type as a sort of general will, a kind of inertia in movement that steams along quite unconsciously in many cases. The vertical descent of will in a particular person is able influence this direction. Some may not keep up or agree to the change of direction, perhaps they drop off and become another race with another destiny. These ideas are connected to the idea of kingship and hierarchy. If the king legitimately decides ``this is the direction'', those who oppose that can not be part of the community. Something to note here is that if we have a goal in mind, the direction has to be changed during the course of travel depending on what is encountered, so it is important to interpret the signs. Nowadays however, we have a situation where there is not much hierarchy, and people become subjected to unhealthy influences instead.

\hfill

\texttt{J. R. R. Tolkien on 2020-09-17 at 08:58 said:}

What else is far from being the product of the evolution of animal species? Tigers? Why do you think they removed the reference to Dacque?

\hfill

\texttt{J. R. R. Tolkien on 2020-09-17 at 09:22 said:}

As we see modern man becoming more and more like an animal, proves that in the beginning man was more and more like a god.

\hfill

\end{sffamily}
\end{footnotesize}

